\documentclass[11pt]{article}

\usepackage[margin=1in]{geometry}
\usepackage{amsmath,amssymb,bm}
\usepackage{booktabs}
\usepackage[hidelinks]{hyperref}

\newcommand{\data}{\mathcal{D}}
\newcommand{\E}{\mathbb{E}}
\newcommand{\Var}{\operatorname{Var}}
\newcommand{\sigmoid}{\operatorname{sigmoid}}

\title{Deriving the p\_free Blend Model as an Approximation to the Population Model}
\author{lenscluster implementation note}
\date{\today}

\begin{document}
\maketitle

\section{Goal}

This note derives the continuous p\_free blend model from the discrete
population model. The derivation is not an identity. It is a controlled
mean-field or plug-in approximation: replace the latent random lens field by
its expectation, and optionally absorb the ignored branch-assignment variance
into the effective scatter.

The conclusion is
\[
    \log p(\data\mid\Theta,\text{blend})
    =
    \ell\!\left(\bar{\bm{\alpha}},\bar{A}\right),
\]
where $\bar{\bm{\alpha}}$ and $\bar{A}$ are the p\_free-weighted mean
deflection and Jacobian fields under the population branch prior.

\section{Population model}

For each independent/free candidate galaxy $i=1,\ldots,N$, introduce a latent
branch assignment
\[
    z_i \in \{0,1\},
    \qquad
    z_i=0 \text{ means scaling branch},\quad
    z_i=1 \text{ means free branch}.
\]
The gate probability is
\[
    q_i
    =
    p(z_i=1)
    =
    p_{\rm free,i}
    =
    \sigmoid\!\left(b_{p(i)}+a\,\eta_i+u_i\right),
\]
where $b_{p(i)}$ is the potfile intercept, $a$ is the magnitude slope,
$\eta_i$ is the standardized magnitude feature, and $u_i$ is the local logit
offset.

Let the scaling and free deflection contributions of candidate $i$ be
\[
    \bm{\alpha}_{S_i}(\bm{x}),
    \qquad
    \bm{\alpha}_{F_i}(\bm{x}),
\]
and define the branch difference
\[
    \Delta\bm{\alpha}_i(\bm{x})
    =
    \bm{\alpha}_{F_i}(\bm{x})-\bm{\alpha}_{S_i}(\bm{x}).
\]
For a branch vector $\bm{z}=(z_1,\ldots,z_N)$, the total deflection field is
\[
    \bm{\alpha}_{\bm{z}}(\bm{x})
    =
    \bm{\alpha}_{\rm other}(\bm{x})
    +
    \sum_{i=1}^N
    \left[
        \bm{\alpha}_{S_i}(\bm{x})
        +
        z_i\,\Delta\bm{\alpha}_i(\bm{x})
    \right].
\]
For local-Jacobian mode, define the analogous Jacobian difference
\[
    \Delta J_i(\bm{x})
    =
    J_{F_i}(\bm{x})-J_{S_i}(\bm{x}),
\]
so that
\[
    A_{\bm{z}}(\bm{x})
    =
    A_{\rm other}(\bm{x})
    +
    \sum_{i=1}^N
    \left[
        J_{S_i}(\bm{x})
        +
        z_i\,\Delta J_i(\bm{x})
    \right].
\]

The exact marginalized population likelihood is
\[
    p(\data\mid\Theta,\bm{q})
    =
    \sum_{\bm{z}\in\{0,1\}^N}
    p(\bm{z}\mid\bm{q})\,
    p(\data\mid\Theta,\bm{z}),
\]
with
\[
    p(\bm{z}\mid\bm{q})
    =
    \prod_{i=1}^N q_i^{z_i}(1-q_i)^{1-z_i}.
\]
In log-likelihood notation,
\[
    p(\data\mid\Theta,\bm{q})
    =
    \E_{\bm{z}\sim{\rm Bernoulli}(\bm{q})}
    \left[
        \exp\left\{
            \ell\!\left(\bm{\alpha}_{\bm{z}},A_{\bm{z}}\right)
        \right\}
    \right].
\]
This expression is exact but expensive because it contains $2^N$ global lens
models.

\section{Mean lens field under the population prior}

Since $z_i$ is Bernoulli with mean $q_i$,
\[
    \E[z_i]=q_i.
\]
Therefore the population-prior mean deflection field is
\[
    \bar{\bm{\alpha}}(\bm{x})
    =
    \E_{\bm{z}}\!\left[\bm{\alpha}_{\bm{z}}(\bm{x})\right]
    =
    \bm{\alpha}_{\rm other}(\bm{x})
    +
    \sum_{i=1}^N
    \left[
        (1-q_i)\bm{\alpha}_{S_i}(\bm{x})
        +
        q_i\bm{\alpha}_{F_i}(\bm{x})
    \right].
\]
Likewise,
\[
    \bar{A}(\bm{x})
    =
    \E_{\bm{z}}\!\left[A_{\bm{z}}(\bm{x})\right]
    =
    A_{\rm other}(\bm{x})
    +
    \sum_{i=1}^N
    \left[
        (1-q_i)J_{S_i}(\bm{x})
        +
        q_iJ_{F_i}(\bm{x})
    \right].
\]
These are exactly the fields used by the blend implementation.

\section{Plug-in approximation}

The simplest approximation is to exchange the nonlinear likelihood expectation
for a likelihood evaluated at the expected field:
\[
    \E_{\bm{z}}
    \left[
        \exp\left\{\ell(\bm{\alpha}_{\bm{z}},A_{\bm{z}})\right\}
    \right]
    \approx
    \exp\left\{
        \ell\!\left(
            \E_{\bm{z}}[\bm{\alpha}_{\bm{z}}],
            \E_{\bm{z}}[A_{\bm{z}}]
        \right)
    \right\}.
\]
Taking the logarithm gives
\[
    \boxed{
    \log p(\data\mid\Theta,\bm{q})
    \approx
    \ell\!\left(\bar{\bm{\alpha}},\bar{A}\right).
    }
\]
This is the continuous blend likelihood. In source-plane mode, the likelihood
uses only $\bar{\bm{\alpha}}$ through the source coordinates
\[
    \bar{\bm{\beta}}_j
    =
    \bm{x}_j-\bar{\bm{\alpha}}(\bm{x}_j).
\]
In local-Jacobian mode, it also uses $\bar{A}_j=\bar{A}(\bm{x}_j)$ to construct
the local covariance.

\section{Delta-method derivation}

The plug-in approximation can be motivated by expanding the log likelihood
around the mean field. Write the random branch perturbation around the mean as
\[
    \delta\bm{\alpha}_{\bm{z}}(\bm{x})
    =
    \bm{\alpha}_{\bm{z}}(\bm{x})-\bar{\bm{\alpha}}(\bm{x})
    =
    \sum_i (z_i-q_i)\Delta\bm{\alpha}_i(\bm{x}),
\]
and analogously
\[
    \delta A_{\bm{z}}(\bm{x})
    =
    \sum_i (z_i-q_i)\Delta J_i(\bm{x}).
\]
By construction,
\[
    \E_{\bm{z}}[\delta\bm{\alpha}_{\bm{z}}]=0,
    \qquad
    \E_{\bm{z}}[\delta A_{\bm{z}}]=0.
\]

Let $h$ denote the vector of all field values used by the likelihood at the
observed image positions:
\[
    h_{\bm{z}}
    =
    \left(\bm{\alpha}_{\bm{z}}(\bm{x}_1),\ldots,\bm{\alpha}_{\bm{z}}(\bm{x}_{N_{\rm img}}),
    A_{\bm{z}}(\bm{x}_1),\ldots,A_{\bm{z}}(\bm{x}_{N_{\rm img}})\right),
\]
and let $\bar{h}=\E[h_{\bm{z}}]$. A second-order Taylor expansion gives
\[
    \ell(h_{\bm{z}})
    \approx
    \ell(\bar{h})
    +
    \bm{g}^{\mathsf T}(h_{\bm{z}}-\bar{h})
    +
    \frac{1}{2}(h_{\bm{z}}-\bar{h})^{\mathsf T}
    H
    (h_{\bm{z}}-\bar{h}),
\]
where $\bm{g}$ and $H$ are the gradient and Hessian of $\ell$ at $\bar{h}$.
Taking the expectation of the log likelihood gives
\[
    \E_{\bm{z}}[\ell(h_{\bm{z}})]
    \approx
    \ell(\bar{h})
    +
    \frac{1}{2}
    \operatorname{tr}\!\left(H\,\Var[h_{\bm{z}}]\right),
\]
because the first-order term has zero mean.

The blend approximation drops the curvature correction:
\[
    \frac{1}{2}
    \operatorname{tr}\!\left(H\,\Var[h_{\bm{z}}]\right)
    \approx 0.
\]
This is reasonable when the scaling and free branches are close, the posterior
keeps $q_i$ near 0 or 1 for most candidates, the likelihood is locally flat
relative to the branch perturbations, or the missing variance is effectively
absorbed by source scatter, scaling scatter, or image-plane scatter.

\section{Gaussian residual example}

For intuition, suppose the core likelihood is a simple Gaussian residual model
\[
    \ell(\bm{r})
    =
    -\frac{1}{2}\bm{r}^{\mathsf T}C^{-1}\bm{r}
    -
    \frac{1}{2}\log |2\pi C|,
\]
and the residual is approximately linear in the branch assignments:
\[
    \bm{r}_{\bm{z}}
    =
    \bar{\bm{r}}
    +
    B(\bm{z}-\bm{q}).
\]
Then
\[
    \E[\bm{r}_{\bm{z}}]=\bar{\bm{r}},
\]
and
\[
    \E[\bm{r}_{\bm{z}}^{\mathsf T}C^{-1}\bm{r}_{\bm{z}}]
    =
    \bar{\bm{r}}^{\mathsf T}C^{-1}\bar{\bm{r}}
    +
    \operatorname{tr}\!\left(C^{-1}B\,\Var[\bm{z}]\,B^{\mathsf T}\right).
\]
Since independent Bernoulli variables have
\[
    \Var[\bm{z}]
    =
    \operatorname{diag}\!\left(q_i(1-q_i)\right),
\]
the exact moment-matched Gaussian would have an additional covariance term
\[
    C_{\rm branch}
    =
    B\,\operatorname{diag}\!\left(q_i(1-q_i)\right)B^{\mathsf T}.
\]
The blend model keeps only the mean residual term and drops this additional
branch-assignment covariance:
\[
    \ell_{\rm blend}
    =
    -\frac{1}{2}\bar{\bm{r}}^{\mathsf T}C^{-1}\bar{\bm{r}}
    -
    \frac{1}{2}\log |2\pi C|.
\]
A more expensive moment-matched approximation would instead use
$C+C_{\rm branch}$. The current blend implementation does not include this
extra covariance explicitly.

\section{Relation to the implemented population approximation}

The exact population model marginalizes all branch assignments jointly:
\[
    p(\data\mid\Theta,\bm{q})
    =
    \sum_{\bm{z}}p(\bm{z}\mid\bm{q})p(\data\mid\Theta,\bm{z}).
\]
The implemented population mode avoids the $2^N$ sum by comparing each
candidate one at a time around the all-scaling model. Let
\[
    \ell_0
    =
    \ell(\text{all candidates scaling}),
\]
and
\[
    \Delta_i
    =
    \ell(\text{candidate }i\text{ free, all others scaling})
    -
    \ell_0.
\]
Then the implemented population approximation is
\[
    \log p(\data\mid\Theta,\bm{q})
    \approx
    \ell_0
    +
    \sum_i
    \log\!\left[(1-q_i)+q_i\exp(\Delta_i)\right].
\]
This approximation keeps $q_i$ as an explicit mixture weight and produces a
posterior free-branch membership diagnostic
\[
    r_i
    =
    \frac{q_i\exp(\Delta_i)}
    {(1-q_i)+q_i\exp(\Delta_i)}.
\]

The blend model makes a different approximation. It marginalizes the branch
state at the level of the lens field before evaluating the likelihood:
\[
    \bm{\alpha}_{\bm{z}}
    \quad\longrightarrow\quad
    \E[\bm{\alpha}_{\bm{z}}]=\bar{\bm{\alpha}},
    \qquad
    A_{\bm{z}}
    \quad\longrightarrow\quad
    \E[A_{\bm{z}}]=\bar{A}.
\]
Thus blend mode is best understood as a deterministic mean-field relaxation of
the latent branch model, not as the same likelihood as the population mode.

\section{When the approximation is good or bad}

The blend approximation is most defensible when:
\begin{itemize}
    \item the scaling and free branch deflections are close for most
    candidates;
    \item most $q_i$ values are near 0 or 1, so $q_i(1-q_i)$ is small;
    \item the source-position likelihood is locally smooth over the branch
    perturbations;
    \item missing branch variance can be absorbed by existing scatter terms.
\end{itemize}

It is least defensible when:
\begin{itemize}
    \item $q_i\approx 0.5$ for many candidates;
    \item the free and scaling branches produce substantially different
    deflections near the observed images;
    \item multiple candidates interact strongly through the same source
    centroid;
    \item the posterior classification of a discrete branch is the main
    scientific target.
\end{itemize}

\section{Computational consequence}

The blend model is fast because it evaluates one mean lens field:
\[
    \Theta \mapsto \bar{\bm{\alpha}},\bar{A}
    \mapsto
    \ell(\bar{\bm{\alpha}},\bar{A}).
\]
The implemented population model is slower because it evaluates one base
likelihood plus one replacement likelihood per candidate:
\[
    \ell_0,\quad \ell_1^F,\ldots,\ell_N^F.
\]
Even when vectorized with \texttt{jax.vmap}, the numerical work still scales
with the number of candidates. In local-Jacobian mode, each replacement also
updates the Jacobian covariance terms, making it substantially slower than
source-plane population mode.

\section{Summary}

\begin{center}
\begin{tabular}{lll}
\toprule
Approximation & Mathematical move & Result\\
\midrule
Exact population &
sum over all $\bm{z}$ &
$\sum_{\bm{z}}p(\bm{z}\mid\bm{q})p(\data\mid\bm{z})$\\
Implemented population &
one-candidate likelihood ratios &
$\ell_0+\sum_i\log[(1-q_i)+q_i e^{\Delta_i}]$\\
Blend &
likelihood at mean field &
$\ell(\E[\bm{\alpha}_{\bm{z}}],\E[A_{\bm{z}}])$\\
\bottomrule
\end{tabular}
\end{center}

The blend model follows from the population model by replacing the random
latent branch lens field with its population-prior mean. The approximation
drops the branch-assignment variance and higher-order interactions. This is
why blend is fast and stable, but why population mode is the closer analogy to
an explicit mixture model.

\end{document}
